\documentclass[10pt, a4paper]{article}
\usepackage[polish]{babel}
\usepackage[utf8]{inputenc}
\usepackage{graphicx} 
\usepackage{amsmath}
\usepackage{amsthm}
\usepackage{amssymb}
\usepackage[utf8]{inputenc}
\usepackage[left=2cm, top=-5cm, right=2cm, bottom=2cm, includeheadfoot, headheight=130pt, a4paper]{geometry}
\usepackage[polish]{babel}
\usepackage[utf8]{inputenc}
\usepackage{polski}
\usepackage[T1]{fontenc}
\usepackage{longtable}
\usepackage{caption}
\usepackage{subcaption}
\usepackage{booktabs}
\usepackage{array}
\usepackage{etoolbox}
\usepackage{xcolor}
\usepackage{listings}

\BeforeBeginEnvironment{verbatim}{\small}
\lstset{language=c++,
                basicstyle=\ttfamily\small,
                keywordstyle=\color{blue}\ttfamily\small,
                stringstyle=\color{red}\ttfamily\small,
                commentstyle=\color{green}\ttfamily\small,
                morecomment=[l][\color{magenta}]{\#},
                extendedchars=true,
                literate={ą}{{\k{a}}}1 {ć}{{\'c}}1 {ę}{{\k{e}}}1 {ł}{{\l{}}}1 {ń}{{\'n}}1 {ó}{{\'o}}1 {ś}{{\'s}}1 {ź}{{\'z}}1 {ż}{{\.z}}1
                         {Ą}{{\k{A}}}1 {Ć}{{\'C}}1 {Ę}{{\k{E}}}1 {Ł}{{\L{}}}1 {Ń}{{\'N}}1 {Ó}{{\'O}}1 {Ś}{{\'S}}1 {Ź}{{\'Z}}1 {Ż}{{\.Z}}1
}

\title{Wybrane Zagadnienia Algebry\\Sprawozdanie 2}
\author{Rafał Włodarczyk 279762}
\date{Czerwiec 2026}

\definecolor{darkyellow}{rgb}{0.8, 0.6, 0.0}

\begin{document}

\maketitle

\section*{Zadanie 4}

\subsection*{a) Struktura przechowująca wielomiany wielu zmiennych rzeczywistych}

Proponowany kontener to mapa \textit{klucz} $\to$ \textit{wartość}, 
w której kluczem jest wektor wykładników (jednomian bez współczynnika), 
a wartością jest współczynnik zadanego jednomianu. Nieobecne w strukturze
klucze (jednomiany) są interpretowane jako mające współczynnik równy zero.
Przykład. Funkcja $f$ (Ćw. 37.) dana wzorem $f=x^3 - x^2\cdot y - x^2 \cdot z$ będzie reprezentowana jako trzy elementy w mapie:

\begin{align*}
\{({\color{red}3}, {\color{green}0}, {\color{blue}0}) \to {\color{darkyellow}1}\}  & \quad (\text{odpowiada } {\color{darkyellow}1} \cdot x^{{\color{red}3}}), \\
\{({\color{red}2}, {\color{green}1}, {\color{blue}0}) \to {\color{darkyellow}-1}\} & \quad (\text{odpowiada } {\color{darkyellow}-1} \cdot x^{{\color{red}2}} \cdot y^{{\color{green}1}}), \\
\{({\color{red}2}, {\color{green}0}, {\color{blue}1}) \to {\color{darkyellow}-1}\} & \quad (\text{odpowiada } {\color{darkyellow}-1} \cdot x^{{\color{red}2}} \cdot z^{{\color{blue}1}}).
\end{align*}

\noindent
Zalety struktury: W przypadku wielomianów zadanych w tym sprawozdaniu (dla których większość współczynników jest równa zero) implementacja z mapą pozwala
uniknąć przechowywania wielu zer, jednocześnie zachowując szybkie odwołania do każdego ze współczynników oraz całej listy $\textit{klucz}$ oraz całej listy $\textit{wartość}$.
Realizacja w C++ wygląda następująco:

\begin{lstlisting}[language=C++]
using monomian_t = std::vector<int64_t>;
using poly_t = std::unordered_map<monomian_t, int64_t, container_hash<monomian_t>>;
\end{lstlisting}

\noindent
Wykorzystujemy kontenery \texttt{vector} oraz $\texttt{unordered\_map}$ z STL oraz \texttt{container\_hash} do haszowania wektorów wykładników z Boost.

\subsection*{b) Porządki na jednomianach}

Mając jednomian $aM = a\cdot x_1^{e_1} \cdot x_2^{e_2} \cdots x_n^{e_n}$, gdzie $a$ jest współczynnikiem, a $e_i$ są wykładnikami zmiennych $x_i$, porządki na jednomianach definiują sposób porównywania dwóch jednomianów w celu ustalenia, który z nich jest "większy" lub "mniejszy"
- zostanie wpierw wykorzystany do porównywania jednomianów w \texttt{PolynomialReduce}.

\subsubsection*{ba) Standard Lex}

Standard Lex porównuje jednomiany zakładając ustaloną kolejność zmiennych (e.g. $x_1 > x_2 > \ldots > x_n$) i porównując wykładniki zmiennych w tej kolejności.

\subsubsection*{bb) Permutowany Lex}

Permutowany Lex jest modyfikacją Standard Lex, która dodatkowo przyjmuje permutację $\pi$. Porównuje jednomiany zakładając ustaloną kolejność zmiennych (e.g. $x_{\pi(1)} > x_{\pi(2)} > \ldots > x_{\pi(n)}$) i porównując wykładniki zmiennych w tej permutowanej kolejności.

\subsubsection*{bc) Graded Lex}

Graded Lex jest hybrydą Standard Lex i porządku stopniowego. Najpierw porównuje jednomiany według sumy wykładników (stopnia), a jeśli są równe to porównuje je według Standard Lex.

\subsection*{c) \texttt{PolynomialReduce}}

\texttt{PolynomialReduce} to algorytm pobierający na wejściu wielomian $f$ oraz ciąg wielomianów $\mathcal{G}=(g_1,\dots,g_n)$. 
Algorytm zwraca ciąg współczynników $(a_1,\dots,a_n)$ oraz wielomian $r$ taki, że:

\begin{align*}
    f = \sum_{i=1}^{n} \alpha_i \cdot g_i + r
\end{align*}

\subsubsection*{Idea algorytmu}

Wykorzystując jeden ze wcześniej zdefiniowanych porządków na jednomianach dokonujemy redukcji wielomianu $f$ względem ciągu $\mathcal{G}$

\begin{enumerate}
    \item Dopóki $f \neq 0$. Znajdź $m=\texttt{LT(f, ord)}$, czyli największy jednomian $m$ w $f$ według porządku \textit{ord}. \\ Ustal $LC(m)=a$ oraz $LM(m)=M$.
    \item Dla każdego $g_i$ w $\mathcal{G}$ znajdź $m_i=\texttt{LT(g\_i, ord)}$, czyli największy jednomian $m_i$ w $g_i$ według \textit{ord}. \\ Ustal $LC(m_i)=b$ oraz $LM(m_i)=N$.
    \item Sprawdź czy $LM(m_i)$ dzieli $LM(m)$. Spełniony musi zostać warunek: $\left(\forall j: LM(m)[j] \geq LM(m_i)[j]\right)$.
    \item Jeśli tak:
    \begin{enumerate}
        \item Wykonaj redukcję: $f := f - \frac{LC(m)}{LC(m_i)} \cdot \frac{LM(m)}{LM(m_i)} \cdot g_i$\
        \item Zaktualizuj współczynnik $\alpha_i := \alpha_i + \frac{LC(m)}{LC(m_i)} \cdot \frac{LM(m)}{LM(m_i)}$.
        \item Wróć do kroku 1.
    \end{enumerate}
    \item Jeśli nie, dodaj $LM(m)$ do $r$(eszty) ze współczynnikiem $LC(m)$ i usuń $m$ z $f$. Wróć do kroku 1.
\end{enumerate}

\subsubsection*{Kod w C++}

\begin{lstlisting}[language=c++]    
inline std::tuple<std::vector<poly_t>, poly_t> 
poly_reduce(
    // Wielomian f
    const poly_t& f, 
    // Ciąg wielomianów G = (g_1, ..., g_n) 
    const std::vector<poly_t>& G, 
    // ord na jednomianach
    bool (*compare)(const monomian_t&, const monomian_t&) 
)
{
    size_t n = G.size();
    std::vector<poly_t> alphas(n); 
    poly_t r;
    poly_t p = f; 

    // (1) Dopóki p != 0
    while (!p.empty()) {
        // LT(f) = aM, LM(f) = M, LC(f) = a
        auto lead_p_it = get_leading_term(p, compare);
        const monomian_t lm_p = lead_p_it->first; 
        int64_t lc_p = lead_p_it->second; 
        
        // Dla każdego g_i w G
        bool divided = false;
        for (size_t i = 0; i < n; ++i) {
            if (G[i].empty()) continue;
            
            // (2) Znajdź LT(g_i) = bN, LM(g_i) = N, LC(g_i) = b
            auto lead_g_it = get_leading_term(G[i], compare);
            const auto& lm_g = lead_g_it->first; 
            int64_t lc_g = lead_g_it->second; 
            
            // (3) Sprawdz czy LM(g_i) dzieli LM(f) (iteruj po wykładnikach)
            bool divisible = true;
            for (size_t j = 0; j < lm_p.size(); ++j) {
                if (lm_p[j] < lm_g[j]) {
                    divisible = false;
                    break;
                }
            }
            
            // (4) Jeśli tak, wykonaj redukcję 
            if (divisible) {
                // Oblicz q = LM(f) / LM(g_i) (iteruj po wykładnikach)
                monomian_t q_mon;
                for (size_t j = 0; j < lm_p.size(); ++j) {
                    q_mon.push_back(lm_p[j] - lm_g[j]);
                }

                // Oblicz q_coef = LC(f) / LC(g_i)
                int64_t q_coef = lc_p / lc_g;
                alphas[i][q_mon] += q_coef;
                
                // Zaktualizuj p := p - q_coef * q_mon * g_i
                for (const auto& [mon_g, coef_g] : G[i]) {
                    monomian_t new_mon;
                    for (size_t j = 0; j < mon_g.size(); ++j) {
                        new_mon.push_back(q_mon[j] + mon_g[j]);
                    }
                    p[new_mon] -= coef_g * q_coef;
                    if (p[new_mon] == 0) p.erase(new_mon);
                }
                
                divided = true;
                break;
            }
        }
        
        // (5) Jeśli nie, dodaj LM(f) do r ze współczynnikiem LC(f) i usuń LT(f) z p
        if (!divided) {
            r[lm_p] += lc_p;
            p.erase(lm_p); 
        }
    }

    // Zwracamy krotkę ({alpha_1, ... alpha_n}, r)
    return {alphas, r}; 
}
\end{lstlisting}

\subsection*{d) Rozwiązanie ćwiczenia 37 z listy zadań}

Rozważamy porządek \textit{GradedLex} oraz wielomiany $f=x^3-x^2y-x^2z$, $g_1=x^2y-z$, $g_2=xy-1$. Redukujemy przy użyciu funkcji \texttt{PolynomialReduce}:

Kodowanie parametrów

\begin{lstlisting}[language=c++]
poly_t poly{
    {{3,0,0},{1}},
    {{2,1,0},{-1}},
    {{2,0,1},{-1}}
};

poly_t g1{
    {{2,1,0},{1}},
    {{0,0,1},{-1}},
};

poly_t g2{
    {{1,1,0},{1}},
    {{0,0,0},{-1}},
};
\end{lstlisting}

\begin{verbatim}
r1=[x^3] - [x^2][z] - [z]
a(g1) = -1
a(g2) = 0
r2=[x^3] - [x^2][z] - [x]
a(g2) = -[x]
a(g1) = 0
\end{verbatim}

Interpretacja Wyniku:

\begin{align*}
    f &= -1 \cdot g_1 + 0 \cdot g_2 + (x^3 - x^2z - z) \\
    f &= -x \cdot g_2 + 0 \cdot g_1 + (x^3 - x^2z - x)
\end{align*}

Reszty są różne, ponieważ \texttt{PolynomialReduce} redukuje po pierwszym spełniającym warunek dzielniku. Wynik zależy więc od kolejności wielomianów w ciągu $\mathcal{G}$.

Czy $r_1 - r_2 \in \langle g_1, g_2 \rangle$ Tak ponieważ różnica reszt jest kombinacją liniową $g_1$ i $g_2$.
\begin{align*}
    r_1 - r_2 &= -1 \cdot g_1 + 0 \cdot g_2 + (x^3 - x^2z - z) - (-x \cdot g_2 + 0 \cdot g_1 + (x^3 - x^2z - x)) \\
    &= -1 \cdot g_1 + x \cdot g_2 + (x^3 - x^2z - z - x^3 + x^2z + x) \\
    &= -1 \cdot g_1 + x \cdot g_2 + (x - z) \in \langle g_1, g_2 \rangle
\end{align*}

\subsection*{e) Wpływ porządków na wynik \texttt{Polynomial Reduce}}

Weźmy $h$ takie że dla danych $(a,b,c,d,e,f) = (2,7,9,7,6,2)$:
\begin{align*}
    h(x,y,z) = x^a y^b - y^c z^d + x^e z^f \in \mathbb{R}[x,y,z] 
    h(x,y,z) = x^2 y^7 - y^9 z^7 + x^6 z^2
\end{align*}

Pokażmy, że istnieje taki ciąg wielomianów $\mathcal{G}=(g_1, g_2, g_3)$, że wynik redukcji $h$ względem $\mathcal{G}$ 
będzie różny dla porządków Standard Lex, Permutowany Lex oraz Graded Lex.



\section*{Zadanie 5}

\subsection*{Emblemat z Japońskiego Paszportu w oparciu o Quadrifolium}

\subsection*{Rozmaitość Algebraiczna dla pow. struktury}


\end{document}